\subsection{\gls{dune} \glspl{fd}} \label{sec:GEOMETRY}
\gls{dune}'s \gls{fd} modules are based in the \gls{surf} to survey a wide portfolio of neutrino physics topics. The site has \SI{4700}{\metersofwater} of overburden which significantly reduces the rate of cosmic backgrounds throughout the detector and allows searches for low-energy astroparticle measurements. The PIP-II accelerator upgrade at the \gls{fermilab} will provide a new neutrino beam directed at \gls{surf} to study long-baseline neutrino oscillations. The first Phase of \gls{dune} will deploy two modules with approximately \SI{14}{\kilo\tonne} of active mass. Two distinct designs (or geometries), are going to be implemented: the \gls{hd} and the \gls{vd} designs (see \cref{fig:DUNE_FAR_DETECTORS}), named after the orientation of their drift fields.

The \gls{hd}'s active volume will extend \SI{14.4}{\meter} along the drift coordinate ($x$), \SI{12}{\meter} along the vertical coordinate $y$, and \SI{58.2}{\meter} along the beam-aligned coordinate ($z$). This module will be separated into four distinct drift volumes alternating three \glspl{apa} and two \glspl{cpa}, which are set to a drift potential of \SI{180}{\kilo\volt}, producing an electric field of \SI{500}{\volt\per\centi\meter}. The \gls{hd}'s \glspl{apa} are equipped with three sets of sensing wires, all with a \SI{5}{\milli\metre} pitch and with induction wires oriented at $\pm37.5$ degrees. The induction wires are positioned closest to the open drift volume, observing the signal as the free electrons transition through the layers. The rear \textit{collection} wires absorb the charge of the drift signal and prevent the distortion of the electric field. This results in three semi-independent \textit{views} of a neutrino interaction that can be combined to create time-dependet 2D representations of charged particle interactions in the detector. In the \gls{hd} module, the two sets of induction wires wrap around the \gls{apa}'s centrally positioned collection wires, allowing the simultaneous observation of two drift volumes. This reduces the overall channel count but requires software intervention to deduce which drift volume an induced signal originates from. The two opposing \gls{apa} sides are instrumented with their own collection plane to aid this deduction.

\begin{figure}[ht!]
    \centering
    \includegraphics[width=.49\linewidth]{Images/Geometries/fd10kt.png}
    \includegraphics[width=.49\linewidth]{Images/Geometries/fdvd10kt.png}
  \caption[width=1\linewidth]{(Left) \gls{dune} \gls{fd} \gls{hd} module. (Right) \gls{dune} \gls{fd} \gls{vd} module.\label{fig:DUNE_FAR_DETECTORS}}
\end{figure}

The \gls{fd} \gls{vd} module simplifies the more traditional \gls{hd} design by dividing the active volume into just two drift volumes of dimensions \qtyproduct{13.5 x 6.5 x 60}{\metre}. This has been made possible by advances in the \gls{lar} purification system that allow for longer drift distances without significant charge loss. The module is divided horizontally by the cathode, denoted simply as \textit{Cathode Plane}. The equivalent of the \gls{cpa} in the \gls{vd} module is called the \gls{crp} and, differently to the thin and free frame-tensioned wires of a traditional \gls{lartpc}, its building block is a perforated \gls{pcb} with flat readout strips (\SI{3.125}{\milli\meter} pitch). Similarly to the case of the free wires in the \gls{hd} design, these perforations allow for the drifting electrons to pass by two induction \gls{pcb}-planes, before being collected. This makes the analysis and simulation results very compatible between these two modules. The coordinate system of the \gls{vd} geometry has been rotated by \SI{90}{\degree} along the $z$-axis.

\begin{table}[ht!]
    \centering
    \gls{dune} \gls{fd} Comparison \\
    \begin{tabular}{lcc}
        \toprule
        Parameter & \gls{hd} Geometry & \gls{vd} Geometry \\
        \midrule
        \gls{lar} Mass & \multicolumn{2}{c}{\SI{17.5}{\kilo\tonne}} \\
        Inner Cryostat Volume & \multicolumn{2}{c}{\qtyproduct{15.1 x 14 x 62}{\meter}} \\
        Active \gls{lar} Mass & \SI{13.5}{\kilo\tonne} & \SI{14}{\kilo\tonne} \\
        Nominal \gls{lar} Mass & \SI{10}{\kilo\tonne} & \SI{12.1}{\kilo\tonne} \\
        Anode Modules & 150 \gls{apa}s & 80 \gls{crp}s \\
        Cathode Modules & 300 \gls{cpa}s & 80 Units \\
        Drift Length & \SI{3.6}{\meter} & \SI{7}{\meter} \\
        Drift Field & \SI{500}{\volt\per\centi\meter} & \SI{450}{\volt\per\centi\meter} \\
        TPC Voltage & \SI{-180}{\kilo\volt} & \SI{-350}{\kilo\volt} \\
        TPC Views & 3-view ($\pm$ 37.5°) & 3-view ($\pm$ 30°) \\
        \bottomrule
    \end{tabular}
    \caption[width=1\linewidth]{Key parameters of the \gls{hd} and \gls{vd} \gls{fd} geometries.\label{tab:GEOMETRY_PARAMETERS}}
\end{table}

\Cref{tab:GEOMETRY_PARAMETERS} summarizes the key parameters of the \gls{hd} and \gls{vd} geometries. Given the nature of the \gls{lartpc} technology, in both geometries, the $y$ and $z$ coordinate reconstruction will follow from the \gls{tpc} readout, together with a time-stamp. For beam events, the remaining $x$ coordinate reconstruction follows from a simple calculation of the beam trigger. Note that for non-beam events (such as astrophysical neutrinos), the interaction vertex along the drift volume cannot be reliably reconstructed. To expand the reconstruction capabilities of the experiment, both detectors incorporate a \gls{pds} which makes use of \gls{xa} modules \cite{Machado_2018} distributed over the \gls{apa} (cathode and membrane planes) for \gls{hd} (\gls{vd}). These provide an additional and instantaneous light signal for any interaction which, provided a correct match from the \gls{tpc} readout is found, can be used to reconstruct the drift coordinate of low-energy interacting particles. This enables background fiducialisation along the drift direction. 
%\gls{dune}'s active material, argon, is also unique in its sensitivity to the $\nu_e$ astrophysical flux. The $\nu_e$~CC interaction on $^{40}$Ar has long been known to have a large scattering cross-section suitable for astroparticle measurements~\cite{raghavan1986inverse}. The PDS should allow for improvements in the low energy physics range, specifically regarding supernova and solar neutrinos ($\sim\SI{10}{\mega\electronvolt}$) and allow for a fast vertex reconstruction and improved calorimetric reconstruction of events above $\SI{5}{\mega\electronvolt}$.